\documentclass[11pt,a4paper]{article}
\usepackage[utf8]{inputenc}
\usepackage[T1]{fontenc}
\usepackage[ngerman]{babel}
\usepackage{helvet}
\renewcommand{\familydefault}{\sfdefault}
\usepackage[left=2cm, right=2cm, top=2cm, bottom=2cm]{geometry}
\usepackage{amsmath,amssymb}
\usepackage{array}
\usepackage{booktabs}
\usepackage{tikz}
\usepackage{pgfplots}
\pgfplotsset{compat=1.18}
\usepackage{fancyhdr}

% Minimal Farben
\definecolor{border}{RGB}{221,221,221}
\definecolor{accent}{RGB}{44,90,160}

\pagestyle{fancy}
\fancyhf{}
\fancyhead[L]{\small Physik Klasse 10E}
\fancyhead[R]{\small Probeklausur Kinematik}
\fancyfoot[C]{\small Seite \thepage}
\renewcommand{\headrulewidth}{0.4pt}

\setlength{\parindent}{0pt}
\newcounter{aufgabe}
\newcommand{\aufgabe}[1]{%
\stepcounter{aufgabe}%
\vspace{0.5cm}%
{\bfseries Aufgabe \theaufgabe}\hfill{\small #1 BE}\\[0.2cm]%
}

\begin{document}

% === KOPFBEREICH ===
\begin{center}
{\Large\bfseries Probeklausur: Kinematik}\\[0.3cm]
{\small Klasse 10E \quad|\quad 45 Minuten \quad|\quad 36 BE}
\end{center}

\vspace{0.3cm}
\begin{tabular}{@{}p{6cm}p{8cm}@{}}
Name: & \hrulefill \\[0.3cm]
Datum: & \hrulefill \\
\end{tabular}

\vspace{0.3cm}
{\color{border}\hrule height 0.5pt}
\vspace{0.2cm}

\begin{center}
\small
\begin{tabular}{|c|c|c|c|c|c|c|}
\hline
\textbf{Aufgabe} & 1 & 2 & 3 & 4 & 5 & \textbf{Summe} \\
\hline
\textbf{max. BE} & 6 & 8 & 8 & 6 & 8 & 36 \\
\hline
\textbf{erreicht} & & & & & & \\
\hline
\end{tabular}
\end{center}

\vspace{0.2cm}
{\color{border}\hrule height 0.5pt}
\vspace{0.5cm}

% === AUFGABEN ===

\aufgabe{6}
\textbf{Grundlagen} (AFB I)

\begin{enumerate}
\item[a)] Rechne um: 90 km/h = \underline{\hspace{2cm}} m/s \hfill (1 BE)

\item[b)] Ein Körper fällt 4 Sekunden frei. Berechne die Fallstrecke. \hfill (2 BE)

\vspace{2cm}

\item[c)] Ergänze die Tabelle: \hfill (3 BE)

\begin{center}
\begin{tabular}{|l|c|c|}
\hline
\textbf{Diagramm} & \textbf{Steigung bedeutet} & \textbf{Fläche bedeutet} \\
\hline
$s(t)$-Diagramm & & --- \\
\hline
$v(t)$-Diagramm & & \\
\hline
\end{tabular}
\end{center}
\end{enumerate}

\vspace{0.5cm}

\aufgabe{8}
\textbf{Bremsvorgang} (AFB II)

Ein Auto fährt mit 72 km/h auf einer Landstraße. Der Fahrer erkennt ein Hindernis und bremst mit $a = -4\,\mathrm{m/s^2}$.

\begin{enumerate}
\item[a)] Rechne die Geschwindigkeit in m/s um. \hfill (1 BE)

\vspace{1.5cm}

\item[b)] Berechne die Bremszeit bis zum Stillstand. \hfill (2 BE)

\vspace{2.5cm}

\item[c)] Berechne den Bremsweg. \hfill (3 BE)

\vspace{3cm}

\item[d)] Das Hindernis ist 60 m entfernt. Reicht der Bremsweg? Begründe. \hfill (2 BE)

\vspace{2cm}
\end{enumerate}

\newpage

\aufgabe{8}
\textbf{Freier Fall und senkrechter Wurf} (AFB II)

Ein Ball wird senkrecht nach oben geworfen mit $v_0 = 20\,\mathrm{m/s}$.

\begin{enumerate}
\item[a)] Berechne die Steigzeit (Zeit bis zum höchsten Punkt). \hfill (2 BE)

\vspace{2.5cm}

\item[b)] Berechne die maximale Steighöhe. \hfill (3 BE)

\vspace{3cm}

\item[c)] Mit welcher Geschwindigkeit kommt der Ball wieder unten an? Begründe kurz. \hfill (3 BE)

\vspace{3cm}
\end{enumerate}

\vspace{0.5cm}

\aufgabe{6}
\textbf{Diagramm-Interpretation} (AFB II)

Das $v(t)$-Diagramm zeigt die Bewegung eines Radfahrers:

\begin{center}
\begin{tikzpicture}
\begin{axis}[
    width=10cm, height=5cm,
    xlabel={$t$ in s},
    ylabel={$v$ in m/s},
    xmin=0, xmax=10,
    ymin=0, ymax=12,
    grid=major,
    axis lines=left,
    xtick={0,2,4,6,8,10},
    ytick={0,2,4,6,8,10,12}
]
\addplot[thick, black] coordinates {(0,0) (4,8) (10,8)};
\end{axis}
\end{tikzpicture}
\end{center}

\begin{enumerate}
\item[a)] Beschreibe die Bewegung in Worten (beide Phasen). \hfill (2 BE)

\vspace{2cm}

\item[b)] Berechne die Beschleunigung in den ersten 4 Sekunden. \hfill (2 BE)

\vspace{2cm}

\item[c)] Berechne den Gesamtweg in den 10 Sekunden. \hfill (2 BE)

\vspace{2.5cm}
\end{enumerate}

\newpage

\aufgabe{8}
\textbf{Luftwiderstand und Transfer} (AFB III)

\begin{enumerate}
\item[a)] Ein Fallschirmspringer springt aus großer Höhe. Erkläre physikalisch, warum er nach dem Öffnen des Fallschirms langsamer fällt. Nenne mindestens zwei relevante Größen. \hfill (3 BE)

\vspace{4cm}

\item[b)] Felix Baumgartner erreichte bei seinem Sprung aus 39 km Höhe Überschallgeschwindigkeit (über 1000 km/h). Am Boden wäre seine Grenzgeschwindigkeit nur etwa 50 m/s gewesen. Erkläre diesen Unterschied. \hfill (3 BE)

\vspace{4cm}

\item[c)] Beurteile folgende Aussage: „Schwere Gegenstände fallen immer schneller als leichte."

Ist diese Aussage richtig oder falsch? Begründe deine Antwort physikalisch für den Fall \textit{mit} und \textit{ohne} Luftwiderstand. \hfill (2 BE)

\vspace{4cm}
\end{enumerate}

\vspace{1cm}
{\color{border}\hrule height 0.5pt}
\vspace{0.3cm}
\begin{center}
\textit{Erlaubte Hilfsmittel: Taschenrechner, Formelsammlung}\\[0.2cm]
\textit{Hinweis: Verwende $g \approx 10\,\mathrm{m/s^2}$}
\end{center}

\end{document}
